\documentclass[a4paper,openany,oneside,12pt]{book}

\usepackage[T1]{fontenc}
\usepackage{textcomp}
\usepackage[spanish]{babel} % espanol
\usepackage[utf8]{inputenc} % acentos sin codigo
\renewcommand{\shorthandsspanish}{}

\usepackage{apacite}
\usepackage{titlesec, blindtext, color}
\definecolor{gray75}{gray}{0.75}
\newcommand{\hsp}{\hspace{20pt}}
\titleformat{\chapter}[hang]{\Huge\bfseries}{\thechapter\hsp\textcolor{gray75}{|}\hsp}{0pt}{\Huge\bfseries}
\usepackage{graphicx} % graficos
\usepackage{setspace}
\usepackage{fancyhdr}
\usepackage{eurosym}
\usepackage{verbatim}
\usepackage{xcolor}
%\usepackage{hyperref}
\usepackage{shadow} %Para usar caja con sombra \shabox{}
\usepackage{lettrine}						% Para hacer letras capitales.
\usepackage{rotating}
\usepackage{wasysym} % Para simbolos adicionales
\usepackage{import}
\usepackage{multirow}

\newcommand{\grad}{\hspace{-2mm}$\phantom{a}^{\circ}$}

\clubpenalty=20000	% para que no haya líneas viudas en el texto.
\widowpenalty=20000 % para que no haya huérfanas

% indica a Latex contabilizar el nivel de gerarquía de texto hasta el cuarto nivel;
\setcounter{secnumdepth}{4}
% indica al comando de generación de tabla de contenido que incluya el cuarto nivel:
\setcounter{tocdepth}{4}

%%% Colores de los links
%\hypersetup{colorlinks=false,
%linkcolor=black,
%urlcolor=blue,
%linkcolor=blue}



%%%%%%%%%%%%%%%%%%%%%%%%%%%%%%%%%%%%%%%%%%%%%%%%%%

\begin{document}
\bibliographystyle{apacite}
%\frontmatter
%	\begin{titlepage}

\newcommand{\HRule}{\rule{\linewidth}{0.5mm}} % Defines a new command for the horizontal lines, change thickness here

\center % Center everything on the page

%----------------------------------------------------------------------------------------
%	LOGO SECTION
%----------------------------------------------------------------------------------------

\includegraphics[width=0.6\textwidth]{./portada/logoUB}\\[1cm] % Include a department/university logo - this will require the graphicx package
 
%----------------------------------------------------------------------------------------
%	HEADING SECTIONS
%----------------------------------------------------------------------------------------

\textsc{\LARGE Trabajo de fin de máster}\\[1cm] % Name of your university/college
\textsc{\Large Máster en Arte Sonoro}\\[0.5cm] % Major heading such as course name
\textsc{\large Facultad de Bellas Artes}\\[0.0cm] % Minor heading such as course title

%----------------------------------------------------------------------------------------
%	TITLE SECTION
%----------------------------------------------------------------------------------------

\HRule \\[0.0cm]
{ \Large \bfseries  \appName\ (\textit{GME Modular Emulator})\\ \vspace{0.4cm} \normalsize Una propuesta de emulación digital del sintetizador\\ EMS Synthi 100 del Gabinete de Música Electroacústica\\de Cuenca}\\[0.4cm] % Title of your document
\HRule \\[0.3cm]
 
%----------------------------------------------------------------------------------------
%	AUTHOR SECTION
%----------------------------------------------------------------------------------------

\emph{Autor:}\\
Carlos Arturo \textsc{Guerra~Parra} % Your name
\\[0.5cm]
\emph{Director:} \\
Dr. José Manuel \textsc{Berenguer Alarcón} % Supervisor's Name
\\[0.7cm]

% If you don't want a supervisor, uncomment the two lines below and remove the section above
%\Large \emph{Author:}\\
%John \textsc{Smith}\\[3cm] % Your name

%----------------------------------------------------------------------------------------
%	DATE SECTION
%----------------------------------------------------------------------------------------

{\large 13 de mayo de 2020} % Date, change the \today to a set date if you want to be precise


 
%----------------------------------------------------------------------------------------

\vfill % Fill the rest of the page with whitespace

{\tiny Con el apoyo y la colaboración de las entidades:}\\
\includegraphics[width=0.1\textwidth]{./portada/logoUCLM}
\includegraphics[width=0.1\textwidth]{./portada/logoBellasartes}
\includegraphics[width=0.1\textwidth]{./portada/logoFuzzyGab}

\end{titlepage}
	\newpage
	\mbox{}
	\thispagestyle{empty} % para que no se numere esta página
%	\pagestyle{fancy}
	


	\tableofcontents
	\listoffigures % indice de figuras
	\addcontentsline{toc}{chapter}{Índice de figuras}
	\listoftables
	\addcontentsline{toc}{chapter}{Índice de cuadros}
	

	

% Aquí van los diversos capítulos:

	\import{capitulos/Chapter_Objetivos/}{Chapter_Objetivos.tex}
	\import{capitulos/Chapter_Descripcion_Aplicacion/}{Chapter_Descripcion_Aplicacion.tex}
%	esto es una cita~\cite{Freud} y otra~\cite{Einstein}

%\bibliography{bibtex/biblio} % database is "biblio.bib" located in a "bibtex" subfolder 






\end{document}